\section{Regulatory framework}\label{sec:regulatory}

\sloppy

The audit-trail requirements an ML pipeline must satisfy come from three regulatory sources of growing weight: SR 11-7 / OCC 2011-12 (United States, 2011), the EU AI Act (Regulation 2024/1689, 2024), and the NIST AI Risk Management Framework (NIST AI 100-1, 2023). They differ in legal force and in granularity but converge on a common set of technical record-keeping requirements that the package described in \Cref{sec:package} aims to satisfy.

\subsection{SR 11-7 / OCC 2011-12}\label{sec:sr117}

\citet{FedSR117} and the companion \citet{OCC2011} are the canonical U.S.\ supervisory guidance on model risk management for federally-regulated banking institutions. The guidance applies to ``models'' broadly defined, encompassing both traditional statistical models and modern machine-learning systems. The relevant requirements for audit trails fall in three sections of the guidance:

\textit{Model documentation} (SR 11-7 §V.G). Banks are expected to document each model's purpose, theory, methodology, assumptions, and limitations, with sufficient detail that a knowledgeable reader unfamiliar with the model can understand its operation. For ML models, this includes the choice of architecture, hyperparameters, training data, and validation results.

\textit{Ongoing monitoring} (SR 11-7 §IV.C, within Model Validation). Banks are expected to verify that models perform as intended over time, with specific attention to changes in input data distributions, model drift, and the continued validity of model assumptions. SR 11-7 explicitly mentions monitoring of input variables for drift relative to the development sample.

\textit{Model inventory} (SR 11-7 §V.F). Each model in production must be in a model inventory that captures (at minimum) the model identity, version, owner, validation status, and use cases. The inventory is the durable record an examiner consults during a supervisory review.

The guidance is principles-based; it does not specify how the documentation, monitoring, or inventory should be technically implemented. In practice, large U.S.\ banks operate model risk management organizations of 50 to several hundred staff who maintain these records, often through a combination of centralized governance platforms (e.g., SAS Model Manager) and bespoke spreadsheet-based tracking.

\subsection{The EU AI Act}\label{sec:euaiact}

\citet{EUAIAct} (Regulation (EU) 2024/1689) entered into force on August 1, 2024 with phased implementation: prohibited AI practices became applicable on February 2, 2025, obligations for general-purpose AI models on August 2, 2025, and most remaining obligations including the high-risk regime on August 2, 2026.

For ``high-risk AI systems'' (defined in Article 6 with reference to Annexes I and III), the Act imposes specific technical record-keeping requirements that are stricter than SR 11-7. Two articles are directly relevant.

\textit{Article 12 -- Record-keeping.} High-risk AI systems must technically allow for the automatic recording of events (logs) over the lifetime of the system. The logs must enable monitoring of operation in line with intended purpose and, in particular, identification of situations which may result in the AI system presenting a risk within the meaning of Article 79.

\textit{Article 19 -- Automatically generated logs.} Providers of high-risk AI systems must keep the logs automatically generated by their AI systems for a period determined by Union or national law, and at minimum for six months after the system is placed on the market. The logs must include, at minimum, the period of each use, the reference database against which input data has been checked, the input data for which the search has led to a match, and the identification of the natural persons involved in the verification of the results.

The volatility-forecasting use case demonstrated in \Cref{sec:demonstration} is not currently classified as ``high-risk'' under Annex III, which lists creditworthiness and access to essential services among financial-sector use cases but not market-data forecasting. The AI Act framing is therefore forward-looking for the present demonstration; we record audit data that would satisfy Article 12 and Article 19 in the event that the use case were brought into scope, or in the event that the firm chooses to apply the same standard voluntarily.

\Cref{tab:article19} maps each Article 19 sub-requirement to the specific \texttt{mr-audit} fields that satisfy it, providing a compliance traceability matrix a CRO can hand to a regulator without re-engineering. The \texttt{cross-domain} examples in the package's \texttt{examples/} directory (a sklearn LogisticRegression credit-scorecard scenario; a statsmodels ARIMA macroeconomic forecaster) demonstrate the same field mapping holds for in-scope Annex III.5(b) credit-scoring use cases as well as for the demonstration's volatility forecaster.

\begin{table}[H]
\centering
\caption{EU AI Act Article 19 sub-requirements mapped to \texttt{mr-audit} schema fields. The Article 19 sub-requirements listed here are the four explicit logging items in the regulation's text; ``identification of natural persons'' is the only one that requires firm-side application policy rather than automatic capture by the package.}\label{tab:article19}
\scriptsize
\setlength{\tabcolsep}{4pt}
\begin{tabular}{@{}p{5.0cm}p{6.5cm}p{3.4cm}@{}}
\toprule
Article 19 sub-requirement & \texttt{mr-audit} field(s) & Captured automatically? \\
\midrule
\multirow{2}{*}{Period of each use} & \texttt{timestamp\_utc} & yes (per record) \\
                                   & \texttt{prediction.predicted\_for\_date} & yes (caller-supplied) \\
\midrule
\multirow{3}{*}{Reference database against} & \texttt{data.input\_data\_source} & yes (caller-supplied) \\
\multirow{3}{*}{which input data has been checked} & \texttt{data.input\_data\_hash} & yes (auto-computed) \\
                                                  & \texttt{data.feature\_names} & yes (from extractor) \\
\midrule
\multirow{2}{*}{Input data for which the search} & \texttt{data.feature\_values\_summary} & yes (from extractor) \\
\multirow{2}{*}{has led to a match} & \texttt{data.input\_data\_hash} & yes (deterministic) \\
\midrule
\multirow{2}{*}{Identification of natural persons} & (firm-side application policy) & no -- requires firm input \\
\multirow{2}{*}{involved in verification of results} & e.g.\ \texttt{model.parameters[''reviewer\_id'']} & yes if firm passes it \\
\midrule
\textit{Beyond Article 19, additional captures:} & & \\
Code provenance & \texttt{code.code\_commit\_hash} & yes (from git) \\
Library provenance & \texttt{code.library\_versions} & yes (from \texttt{sys.modules}) \\
Model identity + parameters & \texttt{model.\{name, version, parameters\}} & yes \\
Compute environment & \texttt{compute.\{python\_version, platform\}} & yes \\
Random state & \texttt{random\_seed} & yes (caller-supplied) \\
Tamper-evidence & \texttt{prev\_record\_hash} chain & yes (auto-computed) \\
\bottomrule
\end{tabular}
\end{table}

The two right-most columns separate concerns honestly: automatic capture (the package's responsibility) from caller-supplied data (the firm's responsibility, e.g.\ logging a reviewer's identifier when a human-in-the-loop is part of the AI system). The package does not attempt to enforce the latter; it provides the schema slot.

\subsection{NIST AI Risk Management Framework}\label{sec:nist}

The NIST AI RMF (\citealp{NIST_AI_100_1}) is a voluntary U.S.\ framework structured around four functions: GOVERN (organizational policies and accountability), MAP (context and risk identification), MEASURE (analysis, assessment, and tracking), and MANAGE (prioritization and response). The MEASURE function explicitly references monitoring for distributional shifts in input data, behavioral drift in model outputs, and traceability of decisions. NIST AI 100-1 is not legally binding, but it is increasingly cited by firms as the de facto operational implementation framework that ``operationalizes'' SR 11-7's principles for ML models.

\subsection{ISO/IEC 42001 management system standard}\label{sec:iso42001}

The fourth framework relevant to ML audit trails is \citet{ISO42001} (published December 2023), the international standard for AI management systems. ISO/IEC 42001 is structured as a Plan-Do-Check-Act management system covering organizational context (Clause 4), leadership (5), planning (6), support (7), operation (8), performance evaluation (9), and improvement (10). The audit-trail-relevant requirements concentrate in three clauses:

\textit{Clause 6.1.2 (AI system risk assessment).} Requires documented risk-assessment processes covering identified risks, treatment plans, and residual-risk acceptance. Audit trails support this by providing the verifiable record against which assessments are revisited.

\textit{Clause 8.3 (AI system impact assessment).} Requires documented impact assessments for AI systems including their effects on individuals, groups, and society. Per-prediction trails enable retrospective impact analysis.

\textit{Clause 9.1 (Monitoring, measurement, analysis, and evaluation).} Requires that the organization determine what needs to be monitored and measured, the methods used, when these are performed, and when the results are analyzed and evaluated. Per-prediction logging is the operational substrate.

Unlike SR 11-7 (US bank-specific) or the EU AI Act (EU-jurisdictional), ISO/IEC 42001 is a voluntary international standard accessible to any organization. Firms pursuing AI Management System certification under ISO/IEC 42001 will increasingly look for tooling that satisfies the relevant clauses automatically. The mapping in \Cref{tab:multi_regulation_matrix} identifies the specific schema fields that satisfy each clause.

\subsection{Synthesis: what an audit trail must capture}\label{sec:requirements}

Across the four regulatory regimes reviewed (US bank model-risk-management via SR 11-7 / OCC 2011-12 in §2.1, EU AI Act in §2.2, NIST AI 100-1 in §2.3, and ISO/IEC 42001 above), an audit trail for an ML inference pipeline should record, at minimum, the per-prediction state in \Cref{tab:requirements}. The simpler three-framework view (with SR 11-7 standing in for OCC 2011-12, which is its companion) appears in \Cref{tab:requirements}; the full clause-by-clause cross-regulation matrix, which separates SR 11-7 and OCC 2011-12 into distinct columns for traceability, appears in \Cref{tab:multi_regulation_matrix}.

\begin{table}[H]
\centering
\caption{Audit-trail field requirements derived from SR 11-7, EU AI Act Articles 12 and 19, and NIST AI 100-1. ``\checkmark'' indicates the framework explicitly references the requirement; blank indicates the requirement is implicit in the broader guidance.}\label{tab:requirements}
\footnotesize
\begin{tabular}{lccc}
\toprule
Field & SR 11-7 & EU AI Act & NIST AI 100-1 \\
\midrule
Model identity (name, version) & \checkmark & \checkmark & \checkmark \\
Model parameters & \checkmark & & \checkmark \\
Model artifact identifier & & \checkmark & \\
Code state (commit, libraries) & & & \checkmark \\
Input data identifier or hash & \checkmark & \checkmark & \checkmark \\
Input feature values or summary & \checkmark & & \checkmark \\
Compute environment & & & \\
Random seed (for stochastic models) & & & \checkmark \\
Prediction value & \checkmark & \checkmark & \checkmark \\
Prediction date / horizon & \checkmark & \checkmark & \checkmark \\
Timestamp of prediction & \checkmark & \checkmark & \checkmark \\
Integrity of the log over time & \checkmark & \checkmark & \\
Replayability (reconstruct prediction) & \checkmark & & \checkmark \\
Drift detection (input + output) & \checkmark & \checkmark & \checkmark \\
\bottomrule
\end{tabular}
\end{table}

\subsubsection{The minimum viable audit trail (MVAT-7)}\label{sec:mvat7}

Of the fourteen rows in \Cref{tab:requirements}, seven appear in at least three of the four regulatory regimes (counting SR 11-7 / OCC 2011-12 as one US bank-MRM regime). These seven form what we term the \emph{Minimum Viable Audit Trail} (MVAT-7), the cross-regulatory floor any deployed ML inference pipeline should be expected to maintain:

\begin{enumerate}[leftmargin=*, label=\textbf{MVAT-\arabic*}]
    \item \textbf{Model identity} (name + version) \hfill --- traceability of which model produced this prediction
    \item \textbf{Input data hash or identifier} \hfill --- traceability of which data went into this prediction
    \item \textbf{Code provenance} (commit hash, library versions) \hfill --- replayability of the computation
    \item \textbf{Prediction value and horizon} \hfill --- what the prediction was and what it predicted
    \item \textbf{Timestamp} \hfill --- when the prediction was made
    \item \textbf{Chain integrity hash} \hfill --- tamper-evidence of the record itself
    \item \textbf{Reviewer or process identifier} (where humans are in the loop) \hfill --- accountability for downstream review
\end{enumerate}

The MVAT-7 is the conservative cross-section: a set any one of the four regimes will fully accept, and the operational target for a firm that wishes to satisfy all four simultaneously. The richer schema documented in \Cref{tab:schema} captures additional fields (compute environment, random seed, prediction-date horizon, feature summary, output-drift baseline) that satisfy specific clauses without being universal. We propose MVAT-7 as the working terminology for ``the minimum that satisfies SR 11-7 / OCC 2011-12, EU AI Act Article 19, NIST AI 100-1 MEASURE-2, and ISO/IEC 42001 Clause 9.1 simultaneously,'' so that future practitioners and standards bodies can refer to a single named target.

\subsubsection{Cross-regulation traceability matrix}\label{sec:cross_regulation}

For an organization operating across multiple jurisdictions or pursuing multi-standard certification, the relevant question is not ``do we satisfy SR 11-7?'' or ``do we satisfy the EU AI Act?'' separately but rather ``do we satisfy all simultaneously, and where is the audit-trail field that demonstrates each?'' \Cref{tab:multi_regulation_matrix} provides the clause-by-clause traceability matrix. The matrix is built from the regulatory texts directly (rather than from any one firm's interpretation) and links each row to the specific \texttt{mr-audit} schema field that satisfies it.

\begin{landscape}
\begin{table}
\centering
\caption{Cross-regulation traceability matrix. Each row is an audit-trail requirement; columns map the requirement to the explicit clause in each framework that references it. ``\checkmark'' = explicit reference; ``\textit{impl.}'' = implicit in broader principles; ``--'' = not addressed. The last column identifies the \texttt{mr-audit} schema field or capability that satisfies the requirement.}\label{tab:multi_regulation_matrix}
\footnotesize
\setlength{\tabcolsep}{5pt}
\begin{tabular}{@{}p{4.0cm}p{2.4cm}p{1.9cm}p{2.8cm}p{2.6cm}p{1.6cm}p{5.0cm}@{}}
\toprule
Required field / capability & SR 11-7 & OCC 2011-12 & EU AI Act & NIST AI 100-1 & ISO/IEC 42001 & \texttt{mr-audit} field \\
\midrule
Model identity (name, version) & §V.F, §V.G & §IV.B & Art.~12, Annex~IV~§2 & MEASURE-2 & §6.1.2, §8.3 & \texttt{model.name}, \texttt{.version} \\
Model parameters & §V.G & §IV.B & \textit{impl.} & MEASURE-2 & §8.3 & \texttt{model.parameters} \\
Model artefact hash & \textit{impl.} & \textit{impl.} & Annex~IV~§2 & MAP-3 & §8.3 & \texttt{model.artifact\_hash} \\
Code provenance (commit, libs) & \textit{impl.} & \textit{impl.} & Art.~12 & MEASURE-2 & §8.3 & \texttt{code.commit\_hash}, \texttt{.library\_versions} \\
Input data identifier / hash & §V.G & §IV.B & Art.~19 & MEASURE-2 & §8.3 & \texttt{data.input\_data\_hash} \\
Input feature values & §V.G & §IV.B & Art.~19 & MEASURE-2 & §8.3 & \texttt{data.feature\_values\_summary} \\
Compute environment & -- & -- & \textit{impl.} & MEASURE-2 & §8.3 & \texttt{compute.*} \\
Random seed & -- & -- & -- & MEASURE-2 & \textit{impl.} & \texttt{random\_seed} \\
Prediction value & §V.G & §IV.B & Art.~19 & MEASURE-2 & §8.3 & \texttt{prediction.value} \\
Prediction horizon / date & §V.G & §IV.B & Art.~19 & MEASURE-2 & §8.3 & \texttt{prediction.horizon}, \texttt{.predicted\_for\_date} \\
Timestamp (UTC) & §V.G & §IV.B & Art.~12 & MEASURE-2 & §9.1 & \texttt{timestamp\_utc} \\
Chain integrity & §V.F & §IV.B & Art.~12 & GOVERN-1 & §8.3 & \texttt{prev\_record\_hash} \\
Replayability & §V.G & §IV.B & \textit{impl.} & MEASURE-2 & §8.3 & \texttt{replay()} engine \\
Input drift detection & §IV.C & §IV.B & \textit{impl.} & MEASURE-2 & §9.1 & \texttt{detect\_drift} \\
Output drift detection & §IV.C & §IV.B & \textit{impl.} & MEASURE-2 & §9.1 & \texttt{detect\_output\_drift} \\
Reviewer / human-in-loop ID & §V.G & §IV.B & Art.~14, Art.~19 & GOVERN-2 & §6.1.2 & firm-supplied via \texttt{model.parameters} \\
GDPR explanation & -- & -- & Art.~86; \citet{GDPR}~Art.~22 & GOVERN-3 & §8.3 & \texttt{explain} CLI subcommand \\
Lifecycle retention & §V.G & §IV.B & Art.~19 (min.~6 mo.) & GOVERN-4 & §7.5 & filesystem persistence \\
\bottomrule
\end{tabular}
\end{table}
\end{landscape}

\Cref{tab:multi_regulation_matrix} is, to our knowledge, the first published matrix that maps all five major ML-audit-trail-relevant frameworks (SR 11-7, OCC 2011-12, EU AI Act, NIST AI 100-1, ISO/IEC 42001) clause-by-clause to a specific implementation reference. The matrix is intended as a citable artefact that compliance officers, model-validation analysts, and standards reviewers can reference directly without re-deriving the cross-walks themselves.

The matrix supports two practical claims. First, \texttt{mr-audit}'s schema satisfies every row of the matrix by default, with the two firm-supplied rows (reviewer ID, lifecycle retention) deferred to the firm's operational policy as is conventional. Second, an organization pursuing simultaneous coverage of multiple regulatory regimes does not need separate audit-trail systems per regulation; a single \texttt{mr-audit}-instrumented pipeline satisfies the cross-section.

The package described in \Cref{sec:package} captures every row of \Cref{tab:requirements} by default, with the exception that input feature values are stored as a configurable summary (mean/std or full-vector at the firm's choice) rather than as the raw tick-level inputs, to keep log sizes tractable.

\subsection{Existing tooling, with honest assessment}\label{sec:existing_tools}

Several open-source and commercial tools partially address the requirements in \Cref{tab:requirements}. None, in our reading, provides a complete solution out of the box:

\begin{itemize}[leftmargin=*]
    \item \textbf{MLflow} (Databricks, OSS): rich experiment tracking and model registry, but designed for the train-time exploration phase. Inference-time logging is supported but lacks immutability, hash-chained integrity, and a regulator-export bundle.
    \item \textbf{Weights \& Biases} (commercial, hosted): excellent UI; data leaves the firm's control; expensive at audit-grade volume; no compliance certification.
    \item \textbf{DVC} (Iterative.ai, OSS): data and code version control. No per-prediction inference logging.
    \item \textbf{Pachyderm} (Hewlett Packard Enterprise, OSS): data versioning and pipeline orchestration. Same gap as DVC for inference logging.
    \item \textbf{AWS SageMaker Model Monitor / Google Vertex AI Model Monitoring}: cloud-locked; not portable; the inference data sits on the cloud provider's infrastructure.
    \item \textbf{IBM AI FactSheets}: documentation templates, not runtime logging.
\end{itemize}

We do not claim our package is uniformly superior to these tools. MLflow has a far richer experiment-management UI, Weights \& Biases has hosted infrastructure, DVC has stronger data versioning. Our claim is narrower: \texttt{mr-audit} is the open-source Python package most narrowly designed for the specific requirements of \Cref{tab:requirements}, with low integration friction and no vendor lock-in.

\Cref{tab:tool_comparison} summarizes the comparison along the audit-trail dimensions that matter for SR 11-7 / EU AI Act compliance. The table is intentionally narrow in scope: it does not score the tools on dimensions where they are clearly stronger than \texttt{mr-audit} (e.g., MLflow's experiment-comparison UI, W\&B's hosted dashboards, DVC's data versioning). Within the audit-trail dimensions, \texttt{mr-audit} is the only entry that satisfies all requirements out of the box.

\providecommand{\halfcheck}{$\bullet$}
\begin{table}[H]
\centering
\footnotesize
\setlength{\tabcolsep}{4pt}
\caption{Audit-trail capability comparison across open-source and commercial tooling. \checkmark{} = native support; \halfcheck{} = partial / requires custom code; -- = not provided. Cells reflect the tools' default behavior as of v2024--2025 documentation review.}\label{tab:tool_comparison}
\begin{tabular}{@{}p{4.6cm}cccccc@{}}
\toprule
Capability & \texttt{mr-audit} & MLflow & W\&B & DVC & \shortstack{SageMaker\\Monitor} & \shortstack{IBM AI\\FactSheets} \\
\midrule
Per-prediction logging & \checkmark & \halfcheck & \halfcheck & -- & \checkmark & -- \\
Hash-chain integrity & \checkmark & -- & -- & -- & -- & -- \\
Replay engine (bit-identical) & \checkmark & -- & -- & -- & -- & -- \\
Input drift detection & \checkmark & -- & \halfcheck & -- & \checkmark & -- \\
Output drift detection & \checkmark & -- & \halfcheck & -- & \checkmark & -- \\
Regulator-export bundle & \checkmark & -- & -- & -- & -- & \halfcheck \\
Code-state provenance (commit, libs) & \checkmark & \halfcheck & \halfcheck & \checkmark & -- & -- \\
Vendor-neutral local storage & \checkmark & \checkmark & -- & \checkmark & -- & \checkmark \\
Open source (MIT/Apache) & \checkmark & \checkmark & -- & \checkmark & -- & \halfcheck \\
\bottomrule
\end{tabular}
\end{table}
